\section{可识别不等于可外推也不等于可决策}
\label{sec:13-Machine-Learning/12-Causal-Inference/06}

一家配送平台在杭州做了个漂亮的随机试验：给骑手加一项调度提示，准时率提高 \(8\) 个百分点，区间窄，各项诊断都干净。产品团队据此全国推开。

三个月后复盘，全国准时率几乎没动。有人怀疑试验有问题，可试验没问题——杭州的效应是真的。问题在别处：参与试验的骑手平均在岗时间更长、熟悉路网，提示对他们是锦上添花；新开城市的骑手更多是兼职新手，提示看不懂也来不及看。更麻烦的是，即使效应在别处也成立，全量推送本身有成本，提示挤占了骑手的注意力，一部分收益被抵消掉了。

三件事被混成了一件：那个数字算得对不对、它在别的人群成不成立、按它行动值不值。它们各有各的失败方式。

\subsection{四道门是串联的，后一道的精度补不了前一道的错}

一个可信的因果结论要连过四道门。第一道是目标定义：处理是什么版本、结果看哪个时间窗、目标人群是谁、要的是 ATE 还是 ATT。第二道是识别：这个量能不能写成观测分布的泛函，靠的是哪几条假设。第三道是估计：有限样本下算得准不准，不确定性从哪来。第四道是外推与决策：结论适用于谁，据此行动是否划算。

\begin{center}
\begin{tikzpicture}[>=stealth,line width=0.7pt]

\foreach \i/\x/\t in {1/0/{定义},2/3.1/{识别},3/6.2/{估计},4/9.3/{外推与决策}}
{
  \draw[rounded corners=3pt] (\x-1.15,-0.5) rectangle (\x+1.15,0.5);
  \node at (\x,0) {\footnotesize \t};
}
\draw[->] (1.15,0) -- (1.95,0);
\draw[->] (4.25,0) -- (5.05,0);
\draw[->] (7.35,0) -- (8.15,0);

\node at (0,-1.05) {\footnotesize 处理版本、时间窗};
\node at (0,-1.5) {\footnotesize 目标人群、估计量};
\node at (3.1,-1.05) {\footnotesize 无混杂、正值性};
\node at (3.1,-1.5) {\footnotesize 一致性、无干扰};
\node at (6.2,-1.05) {\footnotesize 模型设定、方差};
\node at (6.2,-1.5) {\footnotesize 交叉拟合、标准误};
\node at (9.3,-1.05) {\footnotesize 效应修饰、成本};
\node at (9.3,-1.5) {\footnotesize 容量、公平约束};

\draw[->,line width=1pt] (9.3,1.6) -- (9.3,0.65);
\draw[->,line width=1pt] (6.2,1.6) -- (6.2,0.65);
\draw[->,line width=1pt] (3.1,1.6) -- (3.1,0.65);
\draw[->,line width=1pt] (0,1.6) -- (0,0.65);
\node at (4.65,1.95) {\footnotesize 每道门各有自己的失败方式，任何一道关上，后面算得再准也没用};

\end{tikzpicture}
\end{center}

\emph{四道门串联而非并联：置信区间只体现第三道门的成绩，它对前两道和第四道一言不发。}

图里最容易被跳过的是从第三道到第四道那个箭头。多数报告在第三道门就收尾了，一个点估计加一条区间，仿佛剩下的只是执行。开头那家平台的失误正落在这一段：数字属于杭州，行动发生在全国。

\subsection{隐藏混杂在第二道门上，加样本救不了}

若存在未观察的 \(U\) 同时影响 \(X\) 与 \(Y\)，后门路径 \(X\leftarrow U\rightarrow Y\) 就无法被已测协变量阻断。在同一观察机制下扩大样本，只会把那个带混杂的关联估得更精确——方差往下掉，偏差纹丝不动。这是识别层的病，估计层的药治不了。

能用的回应都要从外面引入结构。随机化或准实验设计换来一个已知的分配机制；工具变量、负对照、代理变量各自附带一条不可检验的约束；敏感性分析则退一步，不去证明无偏，而是回答一个可以量化的问题：一个未测混杂要强到什么程度，才足以把当前结论抹平。若很弱的混杂就能抹平，结论脆弱；若需要强到不合情理的关联，结论就更站得住。这不是检验假设，是在测结论对假设失效的耐受度。

\subsection{重叠失败时，模型还在说话，数据已经不说了}

若某些 \(C=c\) 下只有处理者、没有对照者，\(m_0(c)\) 就没有同类样本可学。复杂模型照样吐得出数值，那是结构外推，不是数据提供的反事实比较。

补救的方向是把目标改小：只报共同支持区域内的人群效应，或改用重叠加权。这么做提升了可信度，但同时换掉了估计目标——不能把困难样本删掉之后，还声称估的是原总体的 ATE。诚实的写法是承认目标人群变了，并说清它变成了谁。

\subsection{内部有效不蕴含目标人群有效}

试验可以在参与者身上准确识别效应，却在更老、更病重或资源不同的人群里失效。令 \(S=1\) 表示进入研究。若条件效应

\[
\tau(c)=\E\bigl[Y(1)-Y(0)\given C=c\bigr]
\]

只随已测的 \(C\) 变化，且目标总体在这些 \(C\) 上有支持，那么把研究内的条件效应按目标总体的协变量分布重新加权，就能迁移过去。杭州的骑手在岗时长分布和新城不同，只要在岗时长是被测到的效应修饰变量，重新加权是有依据的。

但若进入研究本身还依赖某个未测的效应修饰变量，迁移就没有依据了。``在独立测试集上复现''也不解决这件事——它验证的是同一数据机制下的稳定性。跨地区、跨时期、跨制度的迁移，需要说清哪些结构机制假定保持不变，而这句话得由领域知识来签字。

\subsection{异质效应不自动等于个体化处方}

条件效应 \(\tau(c)\) 能揭示人群差异，但数据驱动地翻找亚组会同时带来选择偏差和高方差，灵活模型还可能把基线风险的差异误读成处理效应的差异。可用的做法是预先指定亚组，或用诚实分割——一部分数据找亚组，另一部分数据估效应——并连同不确定性一起报，而不是只展示效应最大的那个叶子节点。

从效应到政策还要过成本这一关。若 \(\pi(c)\in\set{0,1}\) 是处理规则，其价值是 \(V(\pi)=\E\bigl[Y(\pi(C))\bigr]\)。即使 \(\tau(c)>0\)，处理成本、副作用、容量上限和公平约束也可能让最优行动不是全员处理。因果效应说的是改变结果的方向和幅度，决策规则还得把所有后果放进同一个效用尺度里比较。这套静态语言在\hyperref[ch:092]{``统计决策与可信预测：模型给出概率之后怎样行动''}一章已经备好——状态、行动、损失、Bayes 风险、容量约束；本章补的是另一半：行动改变的是哪一个结果分布。两者缺一不可，决策理论不能拿观察风险冒充处理效应，因果效应也无法独自判定哪种后果值得付多大代价。

\subsection{时间与干扰会把索引撑大}

纵向决策里，先前的处理会改变后续协变量，后续协变量又影响下一次处理和结果。把这些变量当成普通静态混杂丢进回归，往往同时犯下两个错：控制了中介，还留下新的选择偏差。g-公式、边际结构模型和序贯可忽略性正是为这种处理与混杂互相反馈的结构准备的。

网络、传染和市场里，一个单位的处理会改变另一个单位的结果，无干扰假设失效。这时必须显式定义暴露映射、集群或网络层面的干预，``处理组减对照组''不再对应任何一个稳定的个体效应。

这些扩展不只是把公式写长。它们说的是同一件事：潜在结果的索引必须容纳所有真正改变结果的干预维度，漏掉一维，写下来的量就不是想问的那个量。

\subsection{一份可审查的结论把假设放在明处}

因此报告的结构应当照着那四道门走。先写清估计目标——处理的版本、结果的时间窗、目标人群、效应尺度；再写设计——处理怎么分配、样本怎么进来、时间顺序和缺失机制是什么；然后是识别，逐条交代一致性、无混杂、正值性、无干扰凭什么可信，哪条最弱；接着是估计与不确定性的具体做法；诊断部分给出协变量平衡、权重尾部、共同支持、模型敏感性和负对照的结果；最后交代迁移与决策，说明结论适用于谁、什么环境变化会让它失效、行动还受哪些成本约束。

这么写会让报告显得没那么笃定，那正是它的用处。因果推断不把观察数据变成实验，它做的是更克制的事：把``若改变 \(X\)，\(Y\) 会怎样''拆成明确的目标、看得见的数据和写在纸面上的假设，使得那些数据单独确定不了的部分，不再伪装成回归输出里的一个系数。
